\begin{table}[htp]
	\centering
	\renewcommand{\arraystretch}{1.3}
		\begin{tabular}{ | m{7cm} | m{8cm} | }
			\hline
			\rowcolor[HTML]{5B9BD5} 
			\textcolor{white}{\textbf{Requisito de Información}} & \textcolor{white}{\textbf{Descripción}}\\
			\hline
			
			RI-01: Datos de Entrada & El sistema debe ser capaz de procesar conjuntos de datos tabulares que representen el estado físico inicial de la red, concretamente las potencias activa y reactiva en los diferentes nudos eléctricos \\
			\hline
			RI-02: Datos de Salida & El sistema debe manejar y predecir de forma simultánea múltiples variables continuas correspondientes a la respuesta física de la red. Concretamente, las tensiones nodales, las pérdidas de potencia y la potencia reactiva inyectada por los generadores \\
			\hline
			RI-03: Modelos Predictivos & El sistema debe almacenar de forma automática los mejores modelos encontrados en las diferentes iteraciones para su posible uso en un entorno de producción \\
			\hline
			
		\end{tabular} 
		
		\caption{Requisitos de información del sistema}
	\end{table}